\documentclass[12pt]{article}
\usepackage[utf8]{inputenc}
\usepackage[T1]{fontenc}
\usepackage[margin=1in]{geometry}
\usepackage{setspace}
\usepackage{xcolor}
\usepackage[colorlinks=true, urlcolor=blue!60!black, linkcolor=blue!60!black]{hyperref}
\setlength{\parindent}{0em}
\setlength{\parskip}{0.45em}
\renewcommand{\baselinestretch}{1.0}
\begin{document}
\begin{center}
{\Large \textit{Marbury v. Madison}, 5 U.S. (1 Cranch) 137 (1803)\par}
\vspace{0.3em}
{\normalsize Case Brief --- One-Page Edition\par}
\end{center}

\textbf{Memory Jogger.} Judicial review: courts refuse to enforce unconstitutional statutes, and Congress cannot enlarge the Supreme Court's original jurisdiction (\href{https://supreme.justia.com/cases/federal/us/5/137/}{Justia}).

\textbf{Facts.} Outgoing President Adams appointed William Marbury a District of Columbia justice of the peace for five years. The commission was signed and sealed but undelivered when Jefferson took office, and Secretary of State James Madison withheld it.

\textbf{Procedural History.} Marbury sought mandamus directly in the Supreme Court under Section 13 of the Judiciary Act of 1789; the petition was dismissed for lack of original jurisdiction.

\textbf{Votes.} 4--0: Marshall for Paterson, Chase, and Washington; Cushing and Moore did not participate.

\textbf{Holding.} Article III's grant of original jurisdiction is exclusive, so Section 13 was void insofar as it enlarged it; an act repugnant to the Constitution cannot be applied.

\textbf{Reasoning.} Marbury's right vested on signature, and ministerial duties are judicially enforceable, though political acts are ``only politically examinable.'' A written constitution is ``the fundamental and paramount law of the nation''; letting Congress add to original jurisdiction would make Article III ``mere surplusage''; the judicial oath would be ``worse than solemn mockery'' if judges enforced unconstitutional laws.

\textbf{Examples.} \emph{Same-Side:} original jurisdiction over voters' suits is void. \emph{Opposite-Side:} mandamus to recognize a foreign state is refused. \emph{Fence-Sitter:} departmentalism versus judicial supremacy.

\textbf{Critique.} Van Alstyne faulted Marshall's recusal, merits dicta, and question-begging; Bickel named the countermajoritarian difficulty.

\textbf{Key Quotation.} ``It is emphatically the province and duty of the Judicial Department to say what the law is. Those who apply the rule to particular cases must, of necessity, expound and interpret that rule. If two laws conflict with each other, the Courts must decide on the operation of each.'' --- Chief Justice Marshall, 5 U.S. (1 Cranch) at 177.

\end{document}
